\begin{table}[ht]
\centering
\caption{Fiscal procurement-cost implication of the bounded UTG gap.}
\label{tab:procurement_cost_bound}
\begin{threeparttable}
\begin{tabular}{lr}
\toprule
Component & Value \\
\midrule
Lee midpoint & \BPprocCostBoundUTGmidPct{} \\
Lee lower bound & \BPutgBoundLow{} \\
Lee upper bound & \BPutgBoundHigh{} \\
Admissibility calibration & \BPprocCostBoundAdmissibility{} \\
Annual litigated spending & \BPprocCostBoundAnnualSpend{} \\
Annual procurement-cost implication, midpoint & \BPprocCostBound{} \\
Annual procurement-cost implication, lower & \BPprocCostBoundCIlow{} \\
Annual procurement-cost implication, upper & \BPprocCostBoundCIhigh{} \\
\bottomrule
\end{tabular}
\begin{tablenotes}[flushleft]\footnotesize
\item \textit{Notes:} This calculation is not a full welfare estimate. It applies the bounded price gap to a calibrated admissible share of litigated spending and excludes patient health benefits, search costs, and other welfare components. It should be read together with the sourcing decomposition in Table~\ref{tab:utg_reconciliation}.
\end{tablenotes}
\end{threeparttable}
\end{table}

